\exercisegroup

\begin{enumerate}
\def\labelenumi{\arabic{enumi}.}
\tightlist
\item
  让 \(G\) 设图G满足：证明以下三个命题等价：
\end{enumerate}

\begin{enumerate}
\def\labelenumi{(\alph{enumi})}
\item
  \(\mathbf{G}\) 是一个函数图。
\item
  如果 \(\mathbf{X},\mathbf{Y}\) 是任意两个集合，则
\end{enumerate}

\[
\overline {{{\mathbf {G}}}} (\mathbf {X} \cap \mathbf {Y}) = \overline {{{\mathbf {G}}}} (\mathbf {X}) \cap \overline {{{\mathbf {G}}}} (\mathbf {Y}).
\]

\begin{enumerate}
\def\labelenumi{(\alph{enumi})}
\setcounter{enumi}{2}
\tightlist
\item
  关系 \(\mathbf{X} \cap \mathbf{Y} = \varnothing\) 蕴含
\end{enumerate}

\[
\overline {{\mathbf {G}}} ^ {\mathbf {1}} (\mathbf {X}) \cap \overline {{\mathbf {G}}} ^ {\mathbf {1}} (\mathbf {Y}) = \varnothing .
\]

\begin{enumerate}
\def\labelenumi{\arabic{enumi}.}
\setcounter{enumi}{1}
\tightlist
\item
  设 G 是一个图。证明对于每个集合 X，我们有
\end{enumerate}

\[
\mathrm{G} (\mathrm{X}) = \operatorname{pr} _ {2} (\mathrm{G} \cap (\mathrm{X} \times \operatorname{pr} _ {2} \mathrm{G})) \quad \text { 且 } \quad \mathrm{G} (\mathrm{X}) = \mathrm{G} (\mathrm{X} \cap \operatorname{pr} _ {1} \mathrm{G}).
\]

\begin{enumerate}
\def\labelenumi{\arabic{enumi}.}
\setcounter{enumi}{2}
\tightlist
\item
  设 X, Y, Y', Z 为四个集合。证明
\end{enumerate}

\[
(\mathbf {Y} ^ {\prime} \times \mathbf {Z}) \circ (\mathbf {X} \times \mathbf {Y}) = \varnothing \quad \text { 若 } \quad \mathbf {Y} \cap \mathbf {Y} ^ {\prime} = \varnothing
\]

以及

\[
(\mathbf {Y} ^ {\prime} \times \mathbf {Z}) \circ (\mathbf {X} \times \mathbf {Y}) = \mathbf {X} \times \mathbf {Z} \quad \text { 若 } \quad \mathbf {Y} \cap \mathbf {Y} ^ {\prime} \neq \varnothing .
\]

\begin{enumerate}
\def\labelenumi{\arabic{enumi}.}
\setcounter{enumi}{3}
\tightlist
\item
  让 \((\mathbf{G}_t)_{t\in I}\) 为一族图。证明对每个集合 \(X\) 我们拥有
\end{enumerate}

\[
\left(\bigcup_ {i \in I} G _ {i}\right) \langle X \rangle = \bigcup_ {i \in I} G _ {i} \langle X \rangle ,
\]

并且对每个对象 x，

\[
\left(\bigcap_ {i \in I} G _ {i}\right) \langle \{x \} \rangle = \bigcap_ {i \in I} G _ {i} \langle \{x \} \rangle .
\]

给出两个图 \(\mathbf{G}\) , \(\mathbf{H}\) 和一个集合 \(\mathbf{X}\) 使得

\[
(G \cap H) \langle X \rangle \neq G \langle X \rangle \cap H \langle X \rangle .
\]

\begin{enumerate}
\def\labelenumi{\arabic{enumi}.}
\setcounter{enumi}{4}
\tightlist
\item
  让 \((\mathbf{G}_i)_{i\in \mathbf{I}}\) 为一族图，并设 \(\mathbf{H}\) 为一个图。证明
\end{enumerate}

\[
\left(\bigcup_ {i \in I} G _ {i}\right) \circ H = \bigcup_ {i \in I} (G _ {i} \circ H) \quad \text { 且 } \quad H \circ \left(\bigcup_ {i \in I} G _ {i}\right) = \bigcup_ {i \in I} (H \circ G _ {i}).
\]

\begin{enumerate}
\def\labelenumi{\arabic{enumi}.}
\setcounter{enumi}{5}
\item
  一个图 \(\mathbf{G}\) 是函数性的当且仅当对于每一对图 \(\mathbf{H}, \mathbf{H}'\) 我们拥有 \((\mathbf{H} \cap \mathbf{H}') \circ \mathbf{G} = (\mathbf{H} \circ \mathbf{G}) \cap (\mathbf{H}' \circ \mathbf{G})\) .
\item
  设 G、H、K 为三个图。证明关系
\end{enumerate}

\[
(\mathbf {H} \circ \mathbf {G}) \cap \mathbf {K} \subset (\mathbf {H} \cap (\mathbf {K} \circ \overline {{\mathbf {G}}})) \circ (\mathbf {G} \cap (\overline {{\mathbf {H}}} \circ \mathbf {K})).
\]

\begin{enumerate}
\def\labelenumi{\arabic{enumi}.}
\setcounter{enumi}{7}
\tightlist
\item
  让 \(\Re = (X_{i})_{i\in I}\) 并且 \(\mathfrak{S} = (Y_x)_{x\in K}\) 是集合的两个覆盖 \(E\) .
\end{enumerate}

\begin{enumerate}
\def\labelenumi{(\alph{enumi})}
\item
  证明如果 \(\mathfrak{R}\) 并且 \(\mathfrak{S}\) 是集合的分划 \(\mathbf{E}\) 并且如果 \(\mathfrak{R}\) 比……更细 \(\mathfrak{S}\) ，那么对于每一个 \(x \in K\) 存在 \(\iota \in I\) 使得 \(X_{\iota} \subset Y_{x}\) .
\item
  给出两个覆盖的例子 \(\mathfrak{R}\) , \(\mathfrak{S}\) 使得 \(\mathfrak{R}\) 比……更细 \(\mathfrak{S}\) 但使得 (a) 中所述性质不成立。
\item
  给出两个分划的例子 \(\mathfrak{R}\) , \(\mathfrak{S}\) 的 \(\mathbf{E}\) 使得对于每一个 \(x \in \mathbf{K}\) 存在 \(\iota \in I\) 使得 \(X_{\iota} \subset Y_{x}\) 但使得 \(\mathfrak{R}\) 不是……的加细 \(\mathfrak{S}\) .
\end{enumerate}

